\documentclass{article}
\usepackage[T1]{fontenc}
\usepackage{lmodern}
\usepackage{graphicx}
\usepackage{amsmath,amssymb}
\usepackage[a4paper,margin=2cm]{geometry}
\usepackage[absolute,overlay]{textpos}
\setlength{\TPHorizModule}{1mm}
\setlength{\TPVertModule}{1mm}
\usepackage[indonesian]{babel}
\usepackage{hyperref}
\setlength{\emergencystretch}{2em}
% unit: Halaman 6

\begin{document}

% === Halaman 6 (llm) ===

\begin{itemize}
    \item Sayangnya, algoritma \textit{knapsack} sederhana di atas hanya dapat digunakan untuk enkripsi, tetapi tidak untuk dekripsi.
    \item Misalnya, jika diberikan kriptogram $= 32$, maka tentukan $b_1, b_2, ..., b_6$ sedemikian sehingga
    
    \[ 32 = b_1 + 5b_2 + 6b_3 + 11b_4 + 14b_5 + 20b_6 \eqno{(2)} \]
    
    \item Solusi persamaan (2) ini tidak dapat dipecahkan dalam orde waktu polinomial dengan semakin besarnya $n$ (dengan catatan barisan bobot tidak dalam urutan menaik).
    \item Namun, hal inilah yang dijadikan sebagai kekuatan algoritma \textit{knapsack}.
\end{itemize}

\end{document}
